\item El método de Newton para aproximar una raíz de una ecuación $f(x) = 0$ (véase la sección 2.10) se puede adaptar para aproximar una solución de un sistema de ecuaciones $f(x, y) = 0$ y $g(x, y) = 0$. Las superficies $z = f(x, y)$ y $z = g(x, y)$ se intersectan en una curva que intersecta al plano $xy$ en el punto $(r, s)$, que es la solución del sistema. Si una aproximación inicial $(x_1, y_1)$ está cerca de este punto, entonces los planos tangentes a las superficies en $(x_1, y_1)$ se intersectan en una línea recta que intersecta al plano $xy$ en un punto $(x_2, y_2)$, que debe estar más cercano a $(r, s)$. (Compárese con la figura 2.26). Demuestre que
\[ x_2 = x_1 - \frac{f g_y - f_y g}{f_x g_y - f_y g_x} \qquad \text{y} \qquad y_2 = y_1 - \frac{f_x g - f g_x}{f_x g_y - f_y g_x} \]
en donde $f, g$ y sus derivadas parciales están evaluadas en $(x_1, y_1)$. Si se continúa este procedimiento, se obtienen aproximaciones sucesivas $(x_n, y_n)$.
